%! AP Statistics — Unit 4, Lesson 4.2 — Homework
\documentclass[10pt]{article}
\usepackage{apstats-article}
\usepackage{apstats-boxes}

\begin{document}

\pageheader{Unit 4, Lesson 4.2}{Homework}
\namedateperiod

\vspace{0.1in}

\begin{practicebox}
\small

\textbf{1.} A weather app claims there is a 40\% chance of rain on any given
day this week. Design a simulation to estimate the probability that it rains
on at least 4 of the 7 days.

\begin{enumerate}[label=(\alph*), leftmargin=2em, itemsep=8pt]
  \item Describe your simulation design:
        \begin{itemize}
          \item Chance device: \writeline
          \item How outcomes are assigned: \writeline
          \item What one trial represents: \writeline
          \item What the event of interest is: \writeline
        \end{itemize}
  \item Run 20 trials using random digits below.  (Let 0--3 = rain, 4--9 = no rain.)

        {\small\texttt{7 3 0 5 2 8 9 1 4 6 2 7 1 5 8 3 0 4 9 7 6 2 0 1 5 8 3 2 4 1
        9 5 7 0 6 4 3 8 1 2 5 7 9 0 6 4 2 3 8 1 7 6 5 0 4 2 8 3 1 9
        1 4 7 2 5 8 0 3 6 9 4 1 2 0 5 8 3 7 6 1 2 0 4 8 3 5 9 1 7 6
        0 3 8 5 2 9 4 1 7 6 3 0 5 8 2 1 4 7 9 6 2 3 0 1 8 5 7 4 6 9
        5 1 3 8 0 6 2 7 4 9 3 0 1 6 5 2 8 4 7 9}}

        \vspace{4pt}
        Record the number of ``rain'' days for each set of 7 digits (one trial = 7 digits):

        \vspace{6pt}
        \renewcommand{\arraystretch}{1.6}
        \begin{tabularx}{\linewidth}{@{} *{10}{>{\centering\arraybackslash}X} @{}}
        T1 & T2 & T3 & T4 & T5 & T6 & T7 & T8 & T9 & T10 \\
        \hline
        \blank{0.6cm} & \blank{0.6cm} & \blank{0.6cm} & \blank{0.6cm} & \blank{0.6cm} &
        \blank{0.6cm} & \blank{0.6cm} & \blank{0.6cm} & \blank{0.6cm} & \blank{0.6cm} \\
        \end{tabularx}
        \begin{tabularx}{\linewidth}{@{} *{10}{>{\centering\arraybackslash}X} @{}}
        T11 & T12 & T13 & T14 & T15 & T16 & T17 & T18 & T19 & T20 \\
        \hline
        \blank{0.6cm} & \blank{0.6cm} & \blank{0.6cm} & \blank{0.6cm} & \blank{0.6cm} &
        \blank{0.6cm} & \blank{0.6cm} & \blank{0.6cm} & \blank{0.6cm} & \blank{0.6cm} \\
        \end{tabularx}

  \item How many of your 20 trials had at least 4 rain days? \blank{2cm}
  \item Empirical probability: $\hat{p} = \blank{3cm}$
\end{enumerate}

\vspace{8pt}
\textbf{2.} A multiple-choice quiz has 5 questions, each with 4 answer choices (only one correct).
A student guesses randomly on every question.

\begin{enumerate}[label=(\alph*), leftmargin=2em, itemsep=8pt]
  \item What is the probability of guessing any single question correctly? \blank{2cm}
  \item Design a simulation using a single die (or random digits 1--4) to estimate the
        probability that the student gets \textbf{at least 3 correct}.
        Describe the design: \writeline\writeline
  \item Run 15 trials. Record the number correct per trial:
        \vspace{20pt}
  \item Empirical probability of at least 3 correct: $\hat{p} = \blank{3cm}$
\end{enumerate}

\vspace{8pt}
\textbf{3.} (AP-Style) A simulation is run 500 times to estimate the probability
that a newborn is left-handed. In 65 of the 500 trials the simulated newborn is
left-handed. Which of the following is the best description of the empirical
probability?

\begin{enumerate}[label=(\Alph*), leftmargin=2em, itemsep=4pt]
  \item $500/65 \approx 7.69$
  \item $65/500 = 0.13$
  \item $500/65 = 0.13$
  \item $65 \times 500 = 32{,}500$
\end{enumerate}

\vspace{8pt}
\textbf{4.} \textit{Reflection:} Explain in your own words why the Law of Large Numbers
matters when using simulation to estimate probability. What could go wrong if you based
your estimate on only 5 trials?\\[6pt]
\writeline\writeline\writeline

\vspace{8pt}
\textbf{Extension / Preview:} Use a calculator or app to run 50 random trials of a
coin flip and plot the running relative frequency of heads after each trial. Describe
the shape of the graph. (\textit{Preview: Lesson 4.3 derives the exact probability
rules; today's simulation gives us the intuition for why those rules work.})

\end{practicebox}

\end{document}
